% Realized volatility forecasting -- simple Beamer presentation (single file)
% Figures needed in the same Overleaf project (PDF or PNG):
% fig_rv_timeline, fig_liquidity_ranking, fig_training_curves, fig_test_losses,
% fig_cumulative_qlike, fig_liquidity_qlike, fig_per_stock_gain, fig_forecast_vs_actual
\documentclass[aspectratio=169, 10pt]{beamer}

\usetheme{default}
\usefonttheme{professionalfonts}
\definecolor{Ink}{RGB}{30, 35, 45}
\definecolor{Muted}{RGB}{100, 108, 120}
\setbeamercolor{frametitle}{fg=Ink}
\setbeamercolor{title}{fg=Ink}
\setbeamercolor{structure}{fg=Ink}
\setbeamercolor{normal text}{fg=Ink}
\setbeamerfont{frametitle}{size=\large, series=\bfseries}
\setbeamertemplate{navigation symbols}{}
\setbeamertemplate{itemize items}{\textendash}
\setbeamertemplate{caption}{\raggedright\insertcaption\par}
\setbeamerfont{caption}{size=\scriptsize}
\setbeamercolor{caption name}{fg=Muted}
\setbeamertemplate{footline}{%
    \hfill{\scriptsize\color{Muted}\insertframenumber\,/\,\inserttotalframenumber}\hspace{0.4cm}\vspace{0.25cm}}

\usepackage[T1]{fontenc}
\usepackage{lmodern}
\usepackage{amsmath, amssymb}
\usepackage{booktabs}
\usepackage{graphicx}
\usepackage{adjustbox}
\graphicspath{{./}{figures/}}

\title{Realized Volatility Forecasting in Indian Equities}
\subtitle{HAR-RV, EGARCH and LSTM Models across Liquidity Groups}
\author{Partha Guntur (AU25110011) \and Rajvi Panchal (AU25110019)}
\institute{M.Sc. Quantitative Finance, Ahmedabad University}
\date{November 2026}

\begin{document}

% ------------------------------------------------------------------------------
\begin{frame}[plain]
    \titlepage
\end{frame}

% ------------------------------------------------------------------------------
\begin{frame}{Research Question}
    \textbf{Question.} Can an LSTM forecast next-day realized volatility of Indian stocks better
    than the HAR model, and does the answer depend on liquidity?

    \vspace{0.4cm}
    \textbf{Why it matters}
    \begin{itemize}
        \item Volatility is the main input to option pricing, Value-at-Risk, margins and position sizing.
        \item It is not observed directly: it must be measured from prices and then forecast.
    \end{itemize}

    \vspace{0.3cm}
    \textbf{Two hypotheses}
    \begin{itemize}
        \item \textbf{H1:} a flexible model gains most in illiquid stocks, where a linear model fits worst.
        \item \textbf{H2:} it gains only in liquid stocks, where the volatility measure is cleanest.
    \end{itemize}

    \vspace{0.3cm}
    \textbf{Sample.} 95 NIFTY 100 stocks, 5-minute NSE data. Train 2016--2022, validation 2023--2024,
    test 1 Jan 2025 -- 9 Apr 2026 (29,774 stock-days).
\end{frame}

% ------------------------------------------------------------------------------
\begin{frame}{Theory: Realized Volatility}
    \begin{columns}[T]
        \begin{column}{0.46\textwidth}
            Log price: $dp_t = \mu_t\,dt + \sigma_t\,dW_t$

            \vspace{0.15cm}
            True daily variance (integrated variance):
            $$IV_t = \int_{t-1}^{t} \sigma_s^2\,ds$$
            Realized variance from $M$ intraday returns:
            $$RV_t = \sum_{j=1}^{M} r_{t,j}^2 \;\to\; IV_t \quad (M\to\infty)$$
            \small
            \begin{itemize}
                \item 5-minute returns balance precision against bid--ask noise.
                \item Up to 75 returns per day; overnight gaps excluded.
                \item Annualized volatility $=\sqrt{252\,RV_t}$.
            \end{itemize}
        \end{column}
        \begin{column}{0.52\textwidth}
            \begin{figure}
                \centering
                \includegraphics[width=\textwidth]{fig_rv_timeline}
                \caption{Average intraday volatility of the 95 stocks (5-day average). Volatility clusters: calm spells, sudden spikes (COVID-19, June 2024 election), then slow decay.}
            \end{figure}
        \end{column}
    \end{columns}
\end{frame}

% ------------------------------------------------------------------------------
\begin{frame}{Measuring Liquidity}
    \begin{columns}[T]
        \begin{column}{0.4\textwidth}
            Amihud (2002) illiquidity on 5-minute bars:
            $$ILLIQ_i = \frac{1}{N_i}\sum_{k}\frac{|r_{i,k}|}{\text{Rupee volume}_{i,k}}$$
            \small
            \begin{itemize}
                \item Price move per rupee traded: higher means less liquid.
                \item Measured on 2016--2022 only, so groups use no test data.
                \item Three groups: High (32), Medium (31), Low (32).
                \item Agrees with traded value (Spearman $-0.91$).
            \end{itemize}

        \end{column}
        \begin{column}{0.58\textwidth}
            \begin{figure}
                \centering
                \includegraphics[width=\textwidth]{fig_liquidity_ranking}
                \caption{Illiquidity of the 95 stocks (log scale). Reliance, SBI, Infosys and ICICI Bank are the most liquid. Less liquid stocks are more volatile and less persistent.}
            \end{figure}
            \centering\small
            \begin{adjustbox}{max width=\linewidth}\begin{tabular}{lcc}
                \toprule
                Group & Mean vol & AR(1) of $\ln RV$ \\
                \midrule
                High   & 25.1\% & 0.70 \\
                Medium & 28.5\% & 0.65 \\
                Low    & 33.1\% & 0.65 \\
                \bottomrule
            \end{tabular}\end{adjustbox}
        \end{column}
    \end{columns}
\end{frame}

% ------------------------------------------------------------------------------
\begin{frame}{Theory: The HAR Model}
    \textbf{HAR-RV} (Corsi, 2009), one OLS regression per stock:
    \begin{gather*}
    RV_{t+1} = \beta_0 + \beta_d\,RV_t + \beta_w\,RV_t^{(w)} + \beta_m\,RV_t^{(m)} + \varepsilon_{t+1}\\
    RV_t^{(w)} = \tfrac{1}{5}\textstyle\sum_{k=0}^{4} RV_{t-k}, \qquad RV_t^{(m)} = \tfrac{1}{22}\sum_{k=0}^{21} RV_{t-k}
    \end{gather*}
    \begin{columns}[T]
        \begin{column}{0.5\textwidth}
            \small
            \begin{itemize}
                \item \textbf{Heterogeneous market hypothesis:} daily, weekly and monthly traders react to volatility at different horizons.
                \item The three averages mimic \emph{long memory} with four parameters.
                \item \textbf{log-HAR:} the same regression on $\ln RV$; forecast $\exp(\hat y + s^2/2)$ corrects the log bias.
                \item Logs reduce the pull of a few extreme days.
            \end{itemize}
        \end{column}
        \begin{column}{0.46\textwidth}
            \centering
            \small Average coefficients across stocks (2016--2024)
            \vspace{0.1cm}

            \begin{adjustbox}{max width=\linewidth}\begin{tabular}{lcccc}
                \toprule
                & $\beta_d$ & $\beta_w$ & $\beta_m$ & Sum \\
                \midrule
                HAR-RV  & 0.229 & 0.365 & 0.125 & 0.720 \\
                log-HAR & 0.347 & 0.246 & 0.212 & 0.804 \\
                \bottomrule
            \end{tabular}\end{adjustbox}

            \vspace{0.15cm}
            \raggedright\footnotesize In logs the weights are stable across stocks and persistence is high (0.80); in levels they are pulled around by extreme days.
        \end{column}
    \end{columns}
\end{frame}

% ------------------------------------------------------------------------------
\begin{frame}{Theory: EGARCH}
    \textbf{EGARCH(1,1)} (Nelson, 1991) on daily returns, Student-$t$ errors:
    $$r_t = \mu + \sigma_t z_t, \qquad
    \ln\sigma_t^2 = \omega + \beta\ln\sigma_{t-1}^2 + \alpha\big(|z_{t-1}| - E|z_{t-1}|\big) + \gamma\,z_{t-1}$$
    \begin{columns}[T]
        \begin{column}{0.5\textwidth}
            \small
            \begin{itemize}
                \item Log variance: $\sigma_t^2 > 0$ without restrictions.
                \item $\alpha$: size effect; $\gamma<0$: bad news raises volatility more (leverage).
                \item $\beta$: persistence of volatility.
                \item Parameters estimated on 2016--2022 only, then run forward: no look-ahead.
                \item $\hat\sigma_t$ and $z_t$ are the inputs of EGARCH-LSTM.
            \end{itemize}
        \end{column}
        \begin{column}{0.46\textwidth}
            \centering
            \small Estimates across the 95 stocks
            \vspace{0.1cm}

            \begin{adjustbox}{max width=\linewidth}\begin{tabular}{lcc}
                \toprule
                Parameter & Mean & Median \\
                \midrule
                $\alpha$ (size)        & 0.198 & 0.164 \\
                $\gamma$ (asymmetry)   & $-0.042$ & $-0.044$ \\
                $\beta$ (persistence)  & 0.924 & 0.953 \\
                $\nu$ (tail)           & 4.78 & 4.70 \\
                \bottomrule
            \end{tabular}\end{adjustbox}

            \vspace{0.15cm}
            \raggedright\footnotesize $\gamma<0$ for 92\% of stocks; $\nu\approx 5$ means fat tails.
        \end{column}
    \end{columns}
\end{frame}

% ------------------------------------------------------------------------------
\begin{frame}{Theory: LSTM Network}
    \begin{columns}[T]
        \begin{column}{0.44\textwidth}
            \small
            \begin{align*}
                f_t &= \sigma(W_f[h_{t-1},x_t]+b_f)\\
                i_t &= \sigma(W_i[h_{t-1},x_t]+b_i)\\
                o_t &= \sigma(W_o[h_{t-1},x_t]+b_o)\\
                c_t &= f_t\odot c_{t-1} + i_t\odot\tanh(W_c[h_{t-1},x_t]+b_c)\\
                h_t &= o_t\odot\tanh(c_t)
            \end{align*}
            Forget, input and output gates decide what the memory $c_t$ keeps.

            \vspace{0.15cm}
            \textbf{Design}
            \begin{itemize}
                \item Input: last 22 days. HAR-RV-LSTM: $\ln RV, \ln RV^{(w)}, \ln RV^{(m)}$; EGARCH-LSTM: $\ln\hat\sigma_t, z_t$.
                \item 2 layers (64 units) + stock embedding; one model for all stocks.
                \item Output $\exp(z)>0$; QLIKE loss; 3 seeds averaged.
            \end{itemize}
        \end{column}
        \begin{column}{0.54\textwidth}
            \begin{figure}
                \centering
                \includegraphics[width=\textwidth]{fig_training_curves}
                \caption{Train (dashed) and validation (solid) QLIKE, three seeds. HAR inputs reach a much lower loss than EGARCH inputs within 1--2 epochs; early stopping keeps the best validation epoch.}
            \end{figure}
        \end{column}
    \end{columns}
\end{frame}

% ------------------------------------------------------------------------------
\begin{frame}{Theory: Evaluating Forecasts}
    \begin{columns}[T]
        \begin{column}{0.48\textwidth}
            \textbf{Loss functions}
            $$\text{QLIKE} = \frac{RV}{\hat h} - \ln\frac{RV}{\hat h} - 1$$
            $$\text{MSE} = (RV - \hat h)^2$$
            \small
            \begin{itemize}
                \item RV is a noisy proxy of true variance; QLIKE and MSE still rank forecasts correctly (Patton, 2011).
                \item QLIKE depends on the ratio $RV/\hat h$, so it is not dominated by a few extreme days.
                \item It penalizes under-forecasting more than over-forecasting.
            \end{itemize}
        \end{column}
        \begin{column}{0.48\textwidth}
            \textbf{Diebold--Mariano test}
            $$d_t = \overline{L}_t(A) - \overline{L}_t(B),\qquad t_{DM} = \frac{\bar d}{\widehat{se}_{NW}(\bar d)}$$
            \small
            \begin{itemize}
                \item Daily loss difference averaged across stocks; Newey--West errors for autocorrelation.
                \item Checked with a 5-day block bootstrap.
            \end{itemize}
            \normalsize
            \textbf{Mincer--Zarnowitz}
            $$\ln RV_{t+1} = \alpha + \beta \ln \hat h_{t+1} + e_{t+1}$$
            \small A calibrated forecast has $\beta\approx 1$; $R^2$ = share of variation explained.
        \end{column}
    \end{columns}
\end{frame}

% ------------------------------------------------------------------------------
\begin{frame}{Results: Out-of-Sample Accuracy}
    \begin{figure}
        \centering
        \includegraphics[width=0.85\textwidth, height=0.36\textheight, keepaspectratio]{fig_test_losses}
        \caption{Test losses, 29,774 stock-days (lower is better).}
    \end{figure}
    \vspace{-0.15cm}
    \begin{columns}[T]
        \begin{column}{0.52\textwidth}
            \small
            \begin{adjustbox}{max width=\linewidth}\begin{tabular}{lcccc}
                \toprule
                Model & QLIKE & vs HAR-RV & MZ $R^2$ & Fcst/Actual \\
                \midrule
                HAR-RV      & 0.2286 & ---       & 0.402 & 1.21 \\
                log-HAR     & 0.1965 & $-14.0\%$ & 0.462 & 1.08 \\
                HAR-RV-LSTM & \textbf{0.1935} & $-15.4\%$ & \textbf{0.470} & \textbf{1.06} \\
                EGARCH-LSTM & 0.2709 & $+18.5\%$ & 0.348 & 1.39 \\
                \bottomrule
            \end{tabular}\end{adjustbox}
        \end{column}
        \begin{column}{0.44\textwidth}
            \small
            \begin{itemize}
                \item HAR-RV-LSTM is best on every loss.
                \item 14 of its 15.4 points over HAR-RV come from using logs.
                \item HAR-RV over-forecasts by 21\%; the LSTM by 6\%.
                \item EGARCH-LSTM is worst: daily returns lose intraday information.
            \end{itemize}
        \end{column}
    \end{columns}
\end{frame}

% ------------------------------------------------------------------------------
\begin{frame}{Results: Are the Differences Significant?}
    \begin{figure}
        \centering
        \includegraphics[width=0.9\textwidth, height=0.38\textheight, keepaspectratio]{fig_cumulative_qlike}
        \caption{Cumulative daily QLIKE difference against HAR-RV-LSTM. A rising line means the LSTM keeps gaining.}
    \end{figure}
    \vspace{-0.2cm}
    \begin{columns}[T]
        \begin{column}{0.55\textwidth}
            \small
            \begin{adjustbox}{max width=\linewidth}\begin{tabular}{lccc}
                \toprule
                Comparison & Mean diff. & DM $t$ & Holm $p$ \\
                \midrule
                HAR-RV vs HAR-RV-LSTM      & 0.0351 & 6.18 & $<0.001$ \\
                log-HAR vs HAR-RV-LSTM     & 0.0031 & 2.80 & 0.010 \\
                EGARCH-LSTM vs HAR-RV-LSTM & 0.0776 & 9.89 & $<0.001$ \\
                \bottomrule
            \end{tabular}\end{adjustbox}
        \end{column}
        \begin{column}{0.41\textwidth}
            \small
            \begin{itemize}
                \item All gains are significant; bootstrap intervals exclude zero.
                \item Against log-HAR the gain is small and builds after mid-2025.
                \item Adding EGARCH inputs to HAR-RV-LSTM: $t = 0.24$, no gain.
            \end{itemize}
        \end{column}
    \end{columns}
\end{frame}

% ------------------------------------------------------------------------------
\begin{frame}{Analysis: Does Liquidity Matter?}
    \begin{columns}[T]
        \begin{column}{0.49\textwidth}
            \begin{figure}
                \centering
                \includegraphics[width=\textwidth]{fig_liquidity_qlike}
                \caption{Test QLIKE by liquidity group.}
            \end{figure}
        \end{column}
        \begin{column}{0.49\textwidth}
            \begin{figure}
                \centering
                \includegraphics[width=\textwidth]{fig_per_stock_gain}
                \caption{Per-stock gain of HAR-RV-LSTM over log-HAR.}
            \end{figure}
        \end{column}
    \end{columns}
    \begin{columns}[T]
        \begin{column}{0.44\textwidth}
            \small
            \begin{adjustbox}{max width=\linewidth}\begin{tabular}{lccc}
                \toprule
                Group & log-HAR & HAR-RV-LSTM & Gain \\
                \midrule
                High   & 0.1920 & 0.1915 & 0.3\% \\
                Medium & 0.1925 & 0.1898 & 1.4\% \\
                Low    & 0.2050 & 0.1989 & 3.0\% \\
                \bottomrule
            \end{tabular}\end{adjustbox}
        \end{column}
        \begin{column}{0.52\textwidth}
            \small
            \begin{itemize}
                \item Gain rises as liquidity falls ($t$ = 0.38, 2.43, 3.96): supports \textbf{H1}.
                \item Illiquid stocks have noisier, less persistent RV; the network adapts where a fixed linear rule cannot.
            \end{itemize}
        \end{column}
    \end{columns}
\end{frame}

% ------------------------------------------------------------------------------
\begin{frame}{Analysis: Market Conditions and Stability}
    \begin{columns}[T]
        \begin{column}{0.56\textwidth}
            \begin{figure}
                \centering
                \includegraphics[width=\textwidth]{fig_forecast_vs_actual}
                \caption{Average realized and forecast volatility, test period. HAR-RV sits above realized volatility in calm months.}
            \end{figure}
        \end{column}
        \begin{column}{0.42\textwidth}
            \small
            \begin{adjustbox}{max width=\linewidth}\begin{tabular}{lcc}
                \toprule
                Test QLIKE & log-HAR & HAR-RV-LSTM \\
                \midrule
                Calm days (253)   & 0.1971 & \textbf{0.1936} \\
                Normal days (61)  & 0.1944 & \textbf{0.1929} \\
                High-RV days      & \textbf{0.2810} & 0.2920 \\
                \midrule
                Jan--Apr 2025     & \textbf{0.2010} & 0.2018 \\
                Apr--Aug 2025     & 0.1819 & \textbf{0.1776} \\
                Aug--Dec 2025     & 0.1891 & \textbf{0.1811} \\
                Dec 2025--Apr 2026 & 0.2146 & \textbf{0.2136} \\
                \bottomrule
            \end{tabular}\end{adjustbox}
            \begin{itemize}
                \item The LSTM gains on calm and normal days.
                \item On spike days it reacts slightly slower.
                \item Best in 3 of 4 sub-periods; every seed beats log-HAR.
            \end{itemize}
        \end{column}
    \end{columns}
\end{frame}

% ------------------------------------------------------------------------------
\begin{frame}{Conclusion and Applications}
    \begin{columns}[T]
        \begin{column}{0.48\textwidth}
            \textbf{Conclusion}
            \begin{itemize}
                \item HAR-RV-LSTM gives the most accurate forecasts: $-15.4\%$ QLIKE vs HAR-RV, $-1.6\%$ vs log-HAR.
                \item Most of the improvement comes from logs; the network adds a small, significant edge.
                \item The edge grows as liquidity falls (H1).
                \item EGARCH features add nothing beyond intraday RV.
            \end{itemize}
            \vspace{0.2cm}
            \textbf{Limitations}
            \begin{itemize}
                \item One fairly calm test year; current NIFTY 100 members only.
            \end{itemize}
        \end{column}
        \begin{column}{0.48\textwidth}
            \textbf{Applications}
            \begin{itemize}
                \item \textbf{Options:} ATM straddle $\approx 0.8\,\sigma S\sqrt{T}$; a better $\sigma$ means better pricing and hedging.
                \item \textbf{Risk:} VaR and futures margin $= z_\alpha\,\hat\sigma \times$ position.
                \item \textbf{Portfolios:} size $=$ risk budget$/\hat\sigma$ (volatility targeting).
                \item Liquid large caps: log-HAR is enough; mixed-liquidity books: use the LSTM.
            \end{itemize}
            \vspace{0.2cm}
            \textbf{Future work}
            \begin{itemize}
                \item India VIX, jumps, market RV; test through a crisis.
            \end{itemize}
        \end{column}
    \end{columns}
\end{frame}

\end{document}
